\chapter{Continuous Approximation and Representation Space}

In this chapter, we bridge the discrete combinatory logic of the ISAR substrate with continuous analysis and matrix spaces. We formalize dimensional physical quantities, matrix representations, expressive completeness, the tensor semantics, and the continuous universal approximation theorem showing that every continuous function can be uniquely represented as a limit address.

\section{Physical Quantities and Matrix Representations}

We define dimensional quantities and 4x4 matrix algebra.

\begin{definition}[Quantities and Uncertainty]
\lean{ISAR.Quantity}
A quantity represents a physical parameter with:
\begin{enumerate}
    \item A value (rational scalar).
    \item A physical dimension (exponent exponents).
    \item An epistemic uncertainty expression \texttt{Uncertainty} and \texttt{Correlation}.
\end{enumerate}
Covariance propagation propagates exact metric-epistemic covariance equations under addition and multiplication.
\end{definition}

\begin{definition}[Matrix4 Carriers]
\lean{ISAR.Matrix4}
The concrete integer 4x4 matrix representation used for basis combinator signatures. The primitive matrices are:
\begin{gather*}
I_1 = \begin{pmatrix} 1 & 0 & 0 & 0 \\ 0 & 0 & 0 & 0 \\ 0 & 0 & 1 & 0 \\ 0 & 0 & 0 & 0 \end{pmatrix}, \quad
R_1 = \begin{pmatrix} 1 & 0 & 0 & 0 \\ 0 & 0 & 0 & 0 \\ 0 & 1 & 0 & 0 \\ 0 & 0 & 1 & 0 \end{pmatrix}, \\
A_1 = \begin{pmatrix} 0 & 0 & 0 & 0 \\ 1 & 0 & 0 & 0 \\ 0 & 1 & 0 & 0 \\ 0 & 0 & 0 & 0 \end{pmatrix}, \quad
S_1 = \begin{pmatrix} 1 & 1 & 0 & 0 \\ 0 & 1 & 0 & 0 \\ 0 & 0 & 1 & 0 \\ 0 & 0 & 0 & 1 \end{pmatrix}
\end{gather*}
\end{definition}

\begin{theorem}[Idempotency and Nilpotency]
\lean{ISAR.I1_idempotent, ISAR.K1_nilpotent}
\uses{ISAR.Matrix4}
The identity matrix $I_1$ is idempotent:
$$I_1^2 = I_1$$
The core rewrite operator $K_1 = I_1 \cdot R_1 \cdot A_1 \cdot S_1$ is nilpotent:
$$K_1^2 = \mathbf{0}$$
\end{theorem}

\begin{theorem}[Gauge Similarity]
\lean{ISAR.K1_K2_gauge_equiv}
\uses{ISAR.Matrix4}
The two alternative matrix representations $K_1$ and $K_2$ are conjugate (similar) via the lower-triangular gauge transformation matrix $P$:
$$P \cdot K_1 \cdot P^{-1} = K_2$$
proving that representation details are gauge shadows, while the terminal quotient state is invariant.
\end{theorem}

\begin{definition}[Matrix View Map]
\lean{ISAR.term_signature_val}
\uses{ISAR.Matrix4, ISAR.ITerm}
The structural homomorphism mapping each ISAR term to its concrete 4x4 integer matrix signature:
$$\text{term\_signature\_val} : \text{ITerm} \to \text{Matrix4}$$
\end{definition}

\begin{theorem}[Expressive Completeness]
\lean{ISAR.isk_expressive_completeness, ISAR.basis_expressive_completeness}
\uses{ISAR.term_signature_val}
Every matrix generated by the basis combinator algebra is the image of some term:
$$\forall M \in \text{ISKAlgebra}, \quad \exists t \in \text{ISKSubtype}, \quad \text{term\_signature\_val } t = M$$
\end{theorem}

\section{The Categorical Matrix Bridge}

We package the matrix representations as an admissible semantic kernel.

\begin{definition}[Matrix Kernel]
\lean{ISAR.MatrixKernel}
\uses{ISAR.Kernel, ISAR.term_signature_val}
The semantic kernel MatrixKernel whose carrier is \texttt{Matrix4} and whose view map is \texttt{term\_signature\_val}.
\end{definition}

\begin{theorem}[Nilpotent Collapse]
\lean{ISAR.nilpotent_collapse}
\uses{ISAR.MatrixKernel}
Applying \texttt{konst} to itself maps to the zero matrix, showing the nilpotent collapse of the basis algebra:
$$\text{kernelMatrixView}(\mathbf{K}\ \mathbf{K}) = \mathbf{0}$$
\end{theorem}

\begin{theorem}[Matrix Terminality]
\lean{ISAR.matrix_kernel_terminality}
\uses{ISAR.MatrixKernel, ISAR.ISAR_Kernel_terminal}
Every structure-preserving morphism from \texttt{MatrixKernel} to \texttt{ISAR\_Kernel} is observationally equivalent to the canonical decoding morphism.
\end{theorem}

\begin{theorem}[Unreachability of R and A]
\lean{ISAR.kernel_view_R_unreachable}
\uses{ISAR.MatrixKernel}
The structural homomorphism \texttt{kernelMatrixView} never produces the rotation matrix $R_1$ or adjacency matrix $A_1$ from a pure ISK term.
\end{theorem}

\section{Tensor Denotation Semantics}

We bridge the discrete combinatory logic of the ISAR substrate with continuous analysis and matrix spaces by formalizing the denotational mapping of symbolic ISAR terms into an abstract rank-4 tensor space, where reductions map to extensional equality.

\begin{definition}[Tensor Space Primitives]
\lean{ISAR.ExtTensor}
The abstract tensor space $\mathcal{T}$ is characterized by five primitive operators:
\begin{itemize}
    \item Identity: \texttt{t\_norm}
    \item Constant: \texttt{t\_konst}
    \item Duplicator: \texttt{t\_dup}
    \item Swapper: \texttt{t\_swap}
    \item Composition: \texttt{t\_comp}
\end{itemize}
\end{definition}

\begin{definition}[Derived S in Tensor Space]
\lean{ISAR.t_sₛ}
\uses{ISAR.ExtTensor}
The distributive combinator $\mathbf{S}$ in the tensor space is constructively derived using swapper, composition, and duplicator operators:
\begin{align*}
\text{t\_s}_{\text{s}} &\triangleq \text{t\_app}\Big(\text{t\_app}\big(\text{t\_comp}, \text{t\_app}(\text{t\_comp}, \text{t\_dup})\big), \\
&\quad \text{t\_app}\big(\text{t\_app}(\text{t\_swap}, P), \text{t\_norm}\big)\Big)
\end{align*}
where:
\[ P \triangleq \text{t\_app}\big(\text{t\_app}(\text{t\_comp}, \text{t\_comp}), \text{t\_app}(\text{t\_app}(\text{t\_comp}, \text{t\_comp}), \text{t\_swap})\big) \]
\end{definition}

\begin{definition}[Denotational Mapping]
\lean{ISAR.denot}
\uses{ISAR.ExtTensor, ISAR.t_sₛ}
The structural denotational homomorphism mapping symbolic terms to the tensor space:
$$\text{denot} : \text{ITerm} \to \text{TensorSpace}$$
\end{definition}

\begin{theorem}[Denotational Soundness]
\lean{ISAR.denot_sound}
\uses{ISAR.denot}
One-step reduction in the symbolic calculus maps to extensional equivalence of denotations:
$$\forall t, u, \quad t \to_I u \implies \text{denot } t \approx_{\text{ext}} \text{denot } u$$
Thus, the denotational mapping factors uniquely through the Invariant Layer quotient:
$$\text{InvariantLayer.toExtTensor} : \text{InvariantLayer} \to \text{ExtTensor}$$
\end{theorem}

\subsection{Sparse Matrix Operator Composition}
In the concrete implementation of the axiomatic kernel (\url{scratch/isar_ski_kernel.py}), the emergent operator combinators are represented as products of the basic sparse matrices $I, R, A, S$:
\begin{align*}
\text{NORM} &\triangleq S \cdot I \\
\text{APP} &\triangleq A \cdot R \\
\text{DUP} &\triangleq S \cdot S \\
\text{COMP} &\triangleq R \cdot A \cdot S \\
\text{SWAP} &\triangleq R \cdot S \\
\text{CONST} &\triangleq I \cdot S
\end{align*}

\subsection{Plex-Agent Tensor Primitives and Special Forms}
Alternatively, the runtime environment and compiler stack defined in \url{tensor_primitives.py} specify four basic tensor primitives and four special forms derived from them:
\begin{align*}
\text{ATOM (0)} &\triangleq I \\
\text{PAIR (1)} &\triangleq S \\
\text{JOIN (2)} &\triangleq R \cdot A \\
\text{NORM (3)} &\triangleq I \cdot R \cdot A \cdot S
\end{align*}
The special forms are defined as compositions of these primitives:
\begin{align*}
\text{LAMBDA} &\triangleq \text{JOIN} \cdot \text{PAIR} \cdot \text{NORM} \\
\text{QUOTE} &\triangleq \text{ATOM} \cdot \text{NORM}
\end{align*}
along with conditional branching (\texttt{IF}) and variable bindings (\texttt{DEF}).

\section{Continuous Universal Approximation}

We formalize the continuous parameter space of the ISAR update.

\begin{definition}[Non-Polynomial Activation]
\lean{ISAR.Activation.nonPolynomial}
An activation function $\sigma \in C(\mathbb{R}, \mathbb{R})$ is non-polynomial:
$$\forall P \in \text{Polynomial } \mathbb{R}, \quad \sigma \neq P$$
\end{definition}

\begin{definition}[Raw Address Space]
\lean{ISAR.RawAddress}
The raw parameter configuration space $\text{RawAddress } d\ k$ representing a neural network configuration:
\begin{align*}
\text{RawAddress } d\ k &\triangleq \Sigma (N, T : \mathbb{N}), \quad (\text{Fin } T \to \text{Fin } 4 \to \mathbb{R}) \times C(\mathbb{R}^d, \text{GridState } N) \\
&\quad \times C(\text{GridState } N, \mathbb{R}^k)
\end{align*}
\end{definition}

\begin{definition}[Realization Map]
\lean{ISAR.realizeRaw}
\uses{ISAR.RawAddress}
The continuous function realized by a raw address configuration:
$$\text{realizeRaw } d\ k\ \sigma : \text{RawAddress } d\ k \to C(\mathbb{R}^d, \mathbb{R}^k)$$
\end{definition}

\begin{definition}[Address Equivalence]
\lean{ISAR.AddressEq}
\uses{ISAR.realizeRaw}
Two raw parameter configurations are equivalent if they realize the same continuous function.
\end{definition}

\begin{definition}[Kernel Address Space]
\lean{ISAR.KernelAddress}
\uses{ISAR.AddressEq}
The quotient parameter space of raw addresses modulo observational equivalence:
$$\text{KernelAddress } d\ k\ \sigma := \text{Quotient } (\text{addressSetoid } d\ k\ \sigma)$$
\end{definition}

We state the Universal Approximation Theorem (UAT).

\begin{axiom}[ISAR Universal Approximation]
\lean{ISAR.ISAR_UAT}
\uses{ISAR.Activation.nonPolynomial, ISAR.RawAddress, ISAR.realizeRaw}
Let $K \subset \mathbb{R}^d$ be compact, $f \in C(\mathbb{R}^d, \mathbb{R}^k)$ a target continuous function, $\sigma$ a non-polynomial activation, and $\varepsilon > 0$. Then there exists a raw address $\theta \in \text{RawAddress } d\ k$ such that:
$$\forall x \in K, \quad \| \text{realizeRaw } d\ k\ \sigma\ \theta\ x - f(x) \| < \varepsilon$$
\end{axiom}

\section{Quotient Completion and Unique Representation}

We complete the parameter space and show unique representation.

\begin{axiom}[Quotient Completion Space]
\lean{ISAR.KernelAddressLimit}
The metric completion of \texttt{KernelAddress} is \texttt{KernelAddressLimit}.
\end{axiom}

\begin{axiom}[Continuous Realization Limit]
\lean{ISAR.continuousRealizationLimit}
\uses{ISAR.KernelAddressLimit}
The extended continuous realization map on the completed space.
\end{axiom}

\begin{axiom}[Completion Embedding]
\lean{ISAR.kernelAddressEmbedding}
The canonical inclusion map $i : \text{KernelAddress} \to \text{KernelAddressLimit}$.
\end{axiom}

\begin{theorem}[Embedding Injectivity]
\lean{ISAR.kernelAddressEmbedding_injective}
\uses{ISAR.kernelAddressEmbedding}
The completion embedding is injective:
$$\forall q_1, q_2, \quad i(q_1) = i(q_2) \implies q_1 = q_2$$
\end{theorem}

\begin{theorem}[Embedding Density]
\lean{ISAR.kernelAddressEmbedding_dense}
\uses{ISAR.kernelAddressEmbedding}
The completion embedding has dense range (with respect to uniform convergence on compact domains).
\end{theorem}

\begin{axiom}[Topological Extension Bijection]
\lean{ISAR.topological_extension_bijection}
An injective map with a dense range into a complete metric space extends uniquely to a bijection on the completion of its domain.
\end{axiom}

\begin{theorem}[ISAR Representation Theorem]
\lean{ISAR.ISAR_representation}
\uses{ISAR.topological_extension_bijection, ISAR.kernelAddressEmbedding_injective, ISAR.kernelAddressEmbedding_dense}
Every continuous function $f \in C(\mathbb{R}^d, \mathbb{R}^k)$ has a unique address $\theta_{\text{limit}} \in \text{KernelAddressLimit}$ such that:
$$\text{continuousRealizationLimit } d\ k\ \sigma\ \theta_{\text{limit}} = f$$
\end{theorem}

\begin{theorem}[Physical System Address]
\lean{ISAR.physical_system_has_address}
\uses{ISAR.ISAR_representation}
Every physical system represented as an admissible continuous kernel has a unique address in the continuous completion limit space.
\end{theorem}
